\documentclass[11pt,a4paper]{article}

\usepackage[utf8]{inputenc}
\usepackage[T1]{fontenc}
\usepackage{geometry}
\usepackage{graphicx}
\usepackage{amsmath}
\usepackage{amssymb}
\usepackage{booktabs}
\usepackage{array}
\usepackage{titlesec}
\usepackage{hyperref}

\geometry{margin=0.85in}
\hypersetup{
  colorlinks=true,
  linkcolor=black,
  citecolor=black,
  urlcolor=blue
}

% Compact spacing to keep the report concise.
\titlespacing*{\section}{0pt}{1.4ex plus 0.3ex}{0.8ex}
\titlespacing*{\subsection}{0pt}{1.1ex plus 0.3ex}{0.5ex}
\setlength{\textfloatsep}{8pt plus 2pt minus 2pt}
\setlength{\intextsep}{8pt plus 2pt minus 2pt}
\setlength{\abovedisplayskip}{5pt}
\setlength{\belowdisplayskip}{5pt}

\title{%
  \textbf{House Price Prediction via Tree-Based Models,}\\[4pt]
  \textbf{Coordinate Encoding, and Price-per-sqft Stacking}
}

\author{%
  \begin{tabular}{ll}
    Seungho Han & 20211356 \\
    Moonseong Son        & 20221222 \\
    Sungho Kim           & 20241065 \\
    Jaeyun Lim           & 20251303 \\
  \end{tabular}
}

\date{%
  Team 8 \quad $\cdot$ \quad June 9, 2026
}

\begin{document}

\maketitle

\begin{center}
  Code: \url{https://github.com/pw4ch7re3/house-price-regressor}
\end{center}

\section{Introduction}

House price prediction is a supervised regression task: given property
attributes, estimate a continuous sale price. The \textit{USA House Prices}
dataset used here has three characteristics that shape our approach.
First, several of the most informative features are not natively numeric:
location fields such as \texttt{city} and \texttt{zipcode} carry strong price
signal yet cannot be dropped or naively cast to integers.
Second, the target price is severely right-skewed with a few extreme outliers,
so we remove only the most extreme cases and retain realistic noise.
Third, the feature space is imbalanced---living area (\texttt{sqft\_living})
dominates a tree's feature importance, so a model can score reasonably while
essentially predicting from size alone.

This motivates our central hypothesis: a model that over-relies on the dominant
size feature leaves predictive signal on the table, and richer
representations---through more expressive model structure or better feature
encoding---should recover it. We test this with four models spanning that
spectrum: a Multi-Layer Perceptron (MLP), a Decision Tree (DT), a Random Forest
(RF), and a Multi-Layered Gradient Boosting Decision Tree
(mGBDT)~\cite{feng2018mgbdt}. We further study two domain-aware ideas:
\emph{coordinate encoding}, which converts address fields into geographic
coordinates, and \emph{price-per-sqft stacking}, which blends a direct price
model with a price-per-square-foot model to dampen the influence of size and of
price outliers. Throughout, the naive baseline is a Random Forest on
ordinal-encoded categories (\texttt{rf}/\texttt{cat}), and we ask which
combination of model, encoding, and stacking best improves over it.

\input{Background}
\input{Dataset}
\input{Method}
\input{Result}
\input{Conclusion}

\renewcommand{\refname}{references}
\begin{thebibliography}{9}

\bibitem{feng2018mgbdt}
J. Feng, Y. Yu, and Z.-H. Zhou,
``Multi-Layered Gradient Boosting Decision Trees,''
\textit{Advances in Neural Information Processing Systems (NeurIPS)},
vol.~31, pp.~3551--3561, 2018.

\bibitem{friedman2001greedy}
J. H. Friedman,
``Greedy Function Approximation: A Gradient Boosting Machine,''
\textit{The Annals of Statistics}, vol.~29, no.~5, pp.~1189--1232, 2001.

\bibitem{breiman1984cart}
L. Breiman, J. H. Friedman, R. A. Olshen, and C. J. Stone,
\textit{Classification and Regression Trees}.
Wadsworth, 1984.

\bibitem{chen2016xgboost}
T. Chen and C. Guestrin,
``XGBoost: A Scalable Tree Boosting System,''
\textit{Proc.\ 22nd ACM SIGKDD}, pp.~785--794, 2016.

\bibitem{hastie2009elements}
T. Hastie, R. Tibshirani, and J. Friedman,
\textit{The Elements of Statistical Learning: Data Mining, Inference, and
Prediction}, 2nd ed.
New York: Springer, 2009.

\bibitem{friedman2002stochastic}
J. H. Friedman,
``Stochastic Gradient Boosting,''
\textit{Computational Statistics \& Data Analysis}, vol.~38, no.~4,
pp.~367--378, 2002.

\end{thebibliography}

\end{document}
